% !TeX spellcheck = ru_RU
% !TeX root = vkr.tex

\section{Постановка задачи}
\label{sec:task}

Цель работы~--- реализовать пакет параметрического оценивания для библиотеки pysatl-core.

Для достижения цели поставлены следующие задачи:

\begin{enumerate}[itemsep=0.5em]
    \item Разработать общий интерфейс оценивания, обработку ограничений и диагностику для полных и частично заданных семейств.
    \item Реализовать выборочные статистики и непараметрические точечные оценки.
    \item Адаптировать ранее разработанный прототип к новой архитектуре пакета оценивания.
    \item Реализовать методы минимального контраста, минимального расстояния и оценочных уравнений, включая M-оценивание.
    \item Реализовать оценивание параметров по обычным моментам, \mbox{L-моментам} и усечённым \mbox{L-моментам}.
    \item Реализовать байесовское точечное оценивание на основе сопряжённых априорных распределений.
    \item Реализовать параметрический и непараметрический бутстреп для оценивателей и выборочных статистик.
    \item Провести валидацию пакета и сравнение с SciPy по точности и времени вычислений, подготовить примеры использования.
\end{enumerate}

В работе рассматриваются одномерные выборки независимых одинаково распределённых наблюдений. Семейство распределений для параметрического оценивания задаётся пользователем.
